\documentclass[10pt,a4paper]{amsart}
\usepackage[margin=19mm]{geometry}
\usepackage{amsmath,amssymb,mathtools}
\usepackage[scr=rsfs,cal=euler]{mathalfa}
\usepackage{xfrac}
\usepackage[unicode,hidelinks,bookmarks=false]{hyperref}
\newcommand{\ie}{i.e.\ }
\newcommand{\cf}{cf.\ }
\newcommand{\resp}{respectively\ }
\newcommand{\Subsection}{\subsection}
\providecommand{\nobreakdash}{\nobreak\hbox{-}}
\setlength{\emergencystretch}{2em}
\multlinegap0pt
\usepackage[shortlabels]{enumitem}
\usepackage{amssymb}
\usepackage{multirow}
\usepackage{booktabs}
\usepackage{caption}
\usepackage{bm}
\usepackage{mathtools}
\usepackage{stackengine}
\usepackage{xcolor}
\newtheorem{lemma}{Lemma}
\newcommand{\C}{\mathcal{C}}
\newcommand{\F}{\mathbf{F}}
\newcommand{\FF}{\mathbb{F}}
\newcommand{\f}{\mathbf{f}}
\newcommand{\m}{\mathbf{m}}
\title{Separating subsets from their images}
\author{Marco Barbieri, Maru\v{s}a Lek\v{s}e, Primo\v{z} Poto\v{c}nik, Kamilla Rekv\'enyi}
\date{Source: arXiv:2508.20731v2 (CC BY 4.0)}
\begin{document}
\maketitle

\noindent\emph{Source and licence.} Separating subsets from their images. Version arXiv:2508.20731v2 (CC BY 4.0), distributed by arXiv under the Creative Commons Attribution 4.0 International licence, http://creativecommons.org/licenses/by/4.0/, as recorded in the arXiv licence metadata for this exact version and shown on the abstract page https://arxiv.org/abs/2508.20731v2. Full licence notice: \texttt{licenses/arxiv-2508.20731-CC-BY-4.0.txt}.\par\smallskip

\noindent\textit{Editorial scope.} Counted unit: the article's lemma lem:PSL (source0.tex lines 1806--1826). The article's complete original proof of it is delivered with it (source0.tex lines 1827--1877, one \texttt{proof} environment of 51 lines). No mathematics is written, repaired or re-derived: the statement and the proof are the article's own lines, in the article's own notation, and no retained line is substituted or re-flowed.  No further source-local context is needed: every cross-reference of every retained line resolves inside the two delivered spans.  Nothing was omitted from any retained span, so the package carries no omission record.  No retained line contains a citation, so the wrapper carries no bibliography and no bibliography entry.  No retained line contains a figure environment, an image-inclusion command or drawing code, so no image is delivered and the package declares no asset dependency.  Editorial formatting only, in addition: an AMS \texttt{amsart} preamble that restates the article's own declarations and macros used by the retained lines, namely the environments \texttt{lemma} on separate counters, together with the standard packages those lines need.  The retained environments print in this document as Lemma 1 in order of appearance, whereas the article prints the same items with the numbering of its own full document.  The complete native source of the article as distributed in the arXiv e-print for this exact version is bundled unchanged as \texttt{sources/184046/source0.tex}.
\begin{lemma}\label{lem:PSL}
    Let $G$ be a classical group with linear socle and associated natural module $\FF_q^d$ endowed with a subspace action.
    \begin{enumerate}[$(i)$]
        \item If the domain of the action is the set of $k$-subsets, with $k\le d/2$, then
        \[\m(G) \le \left[ d-\lfloor k/2\rfloor \atop k-\lfloor k/2\rfloor\right]_q .\]
        Consequently, if $\C_k$ is the family of such permutation groups, then
        \[ \lim_n\frac{\f_{\C_k}\left(n\right)}{\sqrt{n}} = \lim_n\frac{\F_{\C_k}\left(n\right)}{\sqrt{n}} = 1 \quad \hbox{ for } k \hbox{ even,}\]
        and
        \[ \lim_n\frac{\F_{\C_k}\left(n\right)}{n^{\frac{k+1}{2k}}} \le 1 \quad \hbox{ for } k \hbox{ odd} .\] 
        \item If the domain of the action is the set of pairs of subspaces complementing each other, one of which of dimension $k\le d/2$, then
        \[\m(G) \le q^{k(d-k)}\left[ d-\lfloor k/2\rfloor \atop k-\lfloor k/2\rfloor\right]_q .\]
        Consequently, if $\C_k$ is the family of such permutation groups, then
        \[ \lim_n\frac{\f_{\C_k}\left(n\right)}{\sqrt{n}} = \lim_n\frac{\F_{\C_k}\left(n\right)}{\sqrt{n}} = 1 \quad \hbox{ for } k \hbox{ even,}\]
        and
        \[ \lim_n\frac{\F_{\C_k}\left(n\right)}{n^{\frac{k+1}{2k}}} \le 1 \quad \hbox{ for } k \hbox{ odd} .\]
        \item If the domain of the action is the set of pairs of subspaces, one of dimension $k < d/2$ and the other of dimension $d-k$, the smaller contained in the larger, then
        \[\m(G) \le \left[ \lceil (d-k)/2 \rceil \atop \lceil (d-3k)/2\rceil \right]_q \left[ \lceil(2d -3k) /2 \rceil \atop \lceil k/2\rceil \right]_q .\]
        Consequently, if $\C_k$ is the family of such permutation groups, then
        \[ \lim_n\frac{\F_{\C_k}\left(n\right)}{n^{\frac{k \left\lceil \frac{d - 3k}{2} \right\rceil + (d-k) \left\lceil \frac{k}{2} \right\rceil}{k(2d-3k)}}} \le 1 .\]
    \end{enumerate}
\end{lemma}

\begin{proof}
    We start by considering the action of $G$ on $k$-subspaces. Let
    \[U = \langle e_1,e_2,\ldots, e_{\lfloor k/2 \rfloor} \rangle ,\]
    and consider the subset
    $$A = \left\{ X \in V^{(k)} \,\middle\vert\,  U \leq X\right\}.$$
    For every $g\in G$, we can choose a subspace $X_g \in V^{(k)}$ such that
    \[ U + U^{g^{-1}} \leq X_g. \]
    (Note that such a subspace exists because $\dim (U + U^{g^{-1}}) \le k$.)
    By construction,
    \[ U^g + U \le X_g^g \in A \cap A^g .\]
    Hence $A\cap A^g$ is nonempty, and $A$ is self-separable.
    Observe that
    \[|A| = \left[ d-\lfloor k/2\rfloor \atop k-\lfloor k/2\rfloor\right]_q = q^{\left(k - \left\lceil\frac{k}{2}\right\rceil \right)(d-k)} + O\left( q^{\left(k - \left\lceil\frac{k}{2}\right\rceil \right)(d-k) -1}\right) ,\]
    and
    \[|V^{(k)}|=\left[ d \atop k \right]_q = q^{k(d-k)} + O\left( q^{k(d-k) -1}\right). \]
    Therefore, we can compute that
    \[ |A| \sim n^{\frac{1}{k} \left\lceil \frac{k}{2}\right\rceil} .\]
    This concludes the proof of~$(i)$.

    Let us now consider the action on pairs of complementary subspaces. Suppose that the smallest one has dimension $k$. Define
    \[ U = \langle e_1,e_2,\ldots, e_{\lfloor k/2 \rfloor} \rangle ,\]
    and consider
    $$A = \left\{ (X,Y) \in V^{(k)}\times V^{(d-k)} \,\middle\vert\,  V= X \oplus Y \hbox{ and } U \leq Y\right\}.$$
    For every $g\in G$, we choose $(X_g,Y_g) \in V^{(k)}\times V^{(d-k)}$ so that
    \[ U + U^{g^{-1}} \le X_g  \quad \hbox{and}  \quad V= X_g \oplus Y_g .\]
    A similar argument as in the first case shows that $A$ is not self-separable, because $(X_g,Y_g)^g \in A \cap A^g$. Moreover,
    \[|A| = q^{k(d-k)} \left[ d-\lfloor k/2\rfloor \atop k-\lfloor k/2\rfloor\right]_q = q^{k(d-k)\left(k - \left\lceil\frac{k}{2}\right\rceil \right)(d-k)} + O\left( q^{k(d-k)\left(k - \left\lceil\frac{k}{2}\right\rceil \right)(d-k) -1}\right) ,\]
    and
    \[|\Omega| = q^{k(d-k)}\left[ d \atop k \right]_q = q^{2k(d-k)} + O\left( q^{2k(d-k) -1}\right). \]
    Therefore,
    \[ |A| \sim n^{\frac{1}{k} \left\lceil \frac{k}{2}\right\rceil} ,\]
    which is enough to conclude~$(ii)$.

    Finally, we need to consider the domain of pairs of subspaces, say $X$ and $Y$, so that $X\in V^{(k)}$, $Y\in V^{(d-k)}$ and $X \le Y$. In this case, we have to set
    \[ U = \langle e_1, \ldots, e_{\lfloor k/2 \rfloor} \rangle\] and
    \[ W = \langle e_1, \ldots, e_{\lfloor (d-k)/2 \rfloor} \rangle.\]
    Considering the set
    \[A = \{ (X,Y) \in V^{(k)}\times V^{(d-k)} \mid X \leq Y,\, U \leq X,\, W \leq Y \},\]
    and using similar strategies as before, we conclude that $A$ is not self-separable.
    In particular, we have
    \begin{align*}
        |A| &= \left[ d- \lfloor (d-k)/2\rfloor \atop d - k -\lfloor (d-k)/2\rfloor \right]_q \left[ d - k - \lfloor k/2\rfloor \atop k -\lfloor k/2\rfloor \right]_q
        \\&= \left[ \lceil (d-k)/2 \rceil \atop \lceil (d-3k)/2\rceil \right]_q \left[ \lceil(2d -3k) /2 \rceil \atop \lceil k/2\rceil \right]_q
        \\&= q^{k \left\lceil \frac{d - 3k}{2} \right\rceil + (d-k) \left\lceil \frac{k}{2} \right\rceil} + O\left(q^{k \left\lceil \frac{d - 3k}{2} \right\rceil + (d-k) \left\lceil \frac{k}{2} \right\rceil - 1}\right)
    \end{align*}
    and
    \[|\Omega| = \left[ d-k \atop k \right]_q\left[ d \atop d-k \right]_q = q^{k(2d-3k)} + O\left( q^{k(2d-3k) -1}\right). \]
    Therefore,
    \[ |A| \sim n^{\frac{k \left\lceil \frac{d - 3k}{2} \right\rceil + (d-k) \left\lceil \frac{k}{2} \right\rceil}{k(2d-3k)}} . \]
    This completes the proof of the missing case~$(iii)$.
\end{proof}

\end{document}
